\section{\texorpdfstring{More tactics: \href{https://lean-lang.org/doc/reference/latest/Tactic-Proofs/Tactic-Reference/}{\texttt{simp}}, \href{https://lean-lang.org/doc/reference/latest/Tactic-Proofs/Tactic-Reference/}{\texttt{constructor}}, \href{https://lean-lang.org/doc/reference/latest/Tactic-Proofs/Tactic-Reference/}{\texttt{cases}}, \href{https://lean-lang.org/doc/reference/latest/Tactic-Proofs/Tactic-Reference/}{\texttt{induction}}, \href{https://lean-lang.org/doc/reference/latest/Tactic-Proofs/Tactic-Reference/}{\texttt{unfold}}}{More tactics: simp, constructor, cases, induction, unfold}}\label{sec:04-tactics:04-more-tactics:1}

\subsection{\texorpdfstring{\texttt{simp}: simplify using known simplification lemmas}{simp: simplify using known simplification lemmas}}\label{sec:04-tactics:04-more-tactics:2}

\begin{lstlisting}
theorem simp_example (n : Nat) : n + 0 = n := by
  simp
\end{lstlisting}

\texttt{simp} automatically searches for known ``simplification'' lemmas and applies them, possibly many at once. This is convenient, but it hides \emph{which} facts were used and \emph{why} the proof works, which is bad for learning. \textbf{This book avoids \texttt{simp} and \texttt{rfl}-as-a-shortcut wherever the point is to understand the proof.} Instead, it uses explicit \texttt{rw} steps that name exactly which equality justifies each step, and \texttt{induction} with a fully spelled-out base case and inductive step. \texttt{simp} should be treated as a tool for later, independent proofs, once what it would have done by hand is already understood.

\begin{quote}
Read more. \hyperref[sec:12-working-efficiently:03-simp:1]{Chapter 12, ``\texttt{simp}, now that you understand what it replaces''} covers when it is the \emph{right} efficient choice, once every step no longer needs to be spelled out.
\end{quote}

\subsection{\texorpdfstring{\texttt{constructor}: build a structure/And/Iff by its constructor}{constructor: build a structure/And/Iff by its constructor}}\label{sec:04-tactics:04-more-tactics:3}

\begin{lstlisting}
theorem and_example (P Q : Prop) (hp : P) (hq : Q) : P (*$\wedge$*) Q := by
  constructor
  (*$\cdot$*) exact hp
  (*$\cdot$*) exact hq
\end{lstlisting}

The \texttt{·} (focus dot) lets you address each remaining goal one at a time.

\begin{mathreading}

\texttt{constructor} invokes the introduction rule of the type of the goal (the \(\wedge\)-intro/\(\vee\)-intro of \hyperref[sec:03-propositions-and-proofs:02-logic-recap:1]{Chapter 3, Section 2}, read more below). For a product/conjunction \(P \wedge Q\) it splits the task into proving each factor separately. This reflects the universal property of the product, ``a map into \(P \times Q\) is a pair of maps,'' so proving \(P\) and proving \(Q\) is enough. More generally, for any \texttt{structure} it reduces the goal to one subgoal per field. In a dual way, \texttt{cases} on \(h : P \vee Q\) below invokes the \emph{elimination} rule of a coproduct (the reading of \(\vee\) as a coproduct in \hyperref[sec:03-propositions-and-proofs:05-and-or-not:1]{Chapter 3, Section 5}). To prove anything from \(P \sqcup Q\) it suffices to prove it from each summand, the case analysis \(\iota_1\)/\(\iota_2\).

\end{mathreading}

\begin{quote}
Read more. \hyperref[sec:03-propositions-and-proofs:02-logic-recap:1]{Chapter 3, Section 2} states the introduction/elimination rules for every connective by name, if ``introduction rule''/``elimination rule'' as general terms are new.
\end{quote}

\subsection{\texorpdfstring{\texttt{cases}: case-split on an inductive value or hypothesis}{cases: case-split on an inductive value or hypothesis}}\label{sec:04-tactics:04-more-tactics:4}

\begin{lstlisting}
theorem or_comm_ex {P Q : Prop} (h : P (*$\vee$*) Q) : Q (*$\vee$*) P := by
  cases h with
  | inl hp => exact Or.inr hp
  | inr hq => exact Or.inl hq
\end{lstlisting}

\subsection{\texorpdfstring{\texttt{induction}: proof by induction on a \texttt{Nat} (or other inductive type)}{induction: proof by induction on a Nat (or other inductive type)}}\label{sec:04-tactics:04-more-tactics:5}

\begin{lstlisting}
theorem add_zero_left (n : Nat) : 0 + n = n := by
  induction n with
  | zero => rfl
  | succ k ih =>
    show 0 + (k + 1) = k + 1
    rw [Nat.add_succ, ih]
\end{lstlisting}

This mirrors the mathematical principle of induction.

\[
P(0) \;\land\; \big(\forall k,\ P(k) \to P(k+1)\big) \;\implies\; \forall n,\ P(n)
\]

\texttt{ih} (induction hypothesis) is exactly \(P(k)\), available to prove \(P(k+1)\).

\subsection{\texorpdfstring{\texttt{unfold}: unfold a definition}{unfold: unfold a definition}}\label{sec:04-tactics:04-more-tactics:6}

\begin{lstlisting}
def isZero (n : Nat) : Prop := n = 0

theorem isZero_zero : isZero 0 := by
  unfold isZero
  -- Goal after unfold: 0 = 0
  rfl
\end{lstlisting}

\begin{mathreading}

\texttt{unfold isZero} replaces the defined name with its definiens, exposing \(\mathrm{isZero}(0) = (0 = 0)\). This uses a \emph{definitional} equality: \(\mathrm{isZero} := (n \mapsto n = 0)\) means the two are interchangeable by definition (like unwinding ``\(n\) is even'' to ``\(\exists k,\ n = 2k\)'' in a proof). So the goal becomes the tautology \(0 = 0\), closed by reflexivity.

\end{mathreading}

\subsection{\texorpdfstring{An aside on transparency: \texttt{def} vs.~\texttt{abbrev} vs.~\texttt{opaque}}{An aside on transparency: def vs.~abbrev vs.~opaque}}\label{sec:04-tactics:04-more-tactics:7}

\texttt{unfold isZero} above had to name \texttt{isZero} explicitly, the same way \texttt{rw} names an equation explicitly, because that is exactly what \texttt{unfold} is for. This is a good moment to name the promise from \hyperref[sec:01-basics:02-def-let-implicit:1]{Chapter 1, Section 2}. \texttt{def} is not the only way to introduce a definition, and the alternatives differ precisely in how much they need to be \emph{told about} like this, versus being seen through automatically.

\begin{itemize}
\tightlist
\item
  \textbf{\texttt{def} (semi-reducible, the default).} Requires an explicit instruction to unfold, \texttt{unfold} by name in a tactic proof, exactly as above. This matches the traceability this book has been deliberately practicing since \protect\hyperlink{simp-simplify-using-known-simplification-lemmas}{\texttt{simp} was first introduced} earlier in this file, nothing unfolds silently, and every step names what justified it.
\item
  \textbf{\texttt{abbrev} (reducible, automatically \texttt{@[reducible, inline]}).} Declares the definition \emph{notational}, not merely equal, so Lean's elaborator and automated machinery treat every occurrence as if the body had been written out directly, with no \texttt{unfold} step to ask for. The clean demonstration of this has to wait for tools that actually \emph{search} through definitions automatically. Typeclass resolution (\hyperref[sec:05-rigor-check:01-structure-vs-class:1]{Chapter 5}) searches at reducible transparency by default, so an \texttt{abbrev} (or a \texttt{def} tagged \texttt{@[reducible]}, as \texttt{isPrime} was in \hyperref[sec:03-propositions-and-proofs:06-quantifiers:1]{Chapter 3}) is visible to that search without being unfolded by hand first, while a plain \texttt{def} generally is not.
\item
  \textbf{\texttt{opaque}.} The opposite extreme, not unfoldable at all, by \texttt{unfold} or anything else, even though it still has a definition somewhere. Useful for genuinely hiding an implementation and exposing only the properties proved about it, the Lean equivalent of citing a chosen but unspecified witness (``let \(c\) be \emph{some} element of the nonempty set \(S\)'') and reasoning only from what is known about it, never from how it was built.
\end{itemize}

\begin{mathreading}

This maps onto a distinction already made informally, if rarely by name, in ordinary mathematical writing. A \texttt{def} is a definition cited by name in a proof, ``by definition of \(f\), \ldots{}'', an explicit, visible step. An \texttt{abbrev} is pure notation, the kind silently unfolded without comment, writing \(2n\) in place of ever having introduced a name \(f(n) := 2n\) at all. An \texttt{opaque} value is an axiomatized or black-boxed object, reasoned about only through its stated properties, never its construction. Lean's three transparency levels are exactly these three mathematical habits, made explicit and machine-checked instead of left to convention.

\end{mathreading}

\hyperref[sec:04-tactics:03-reading-failures:1]{← Reading failures} \textbar{} \hyperref[sec:04-tactics:00-index:1]{Index} \textbar{} \hyperref[sec:04-tactics:05-worked-example:1]{Next: Worked example →}
